% !TeX root = ../papers/nc-p3.tex
\section{六次中心商の証明を整な傾きへ降ろす}
\label{sec:rank-fifteen-integral-slope}
\begin{theorem}\label{thm:rank-fifteen-integral-slope}
標数零の通常な階数15窓に付値を取り、その剰余標数を零または奇素数とする。
$\nu(y)>0$ かつ $D=x^3/y$ が整であるとする。
このとき各始点の未切断六次経路は84次元で、
四つの中心四次式を選んだ商の六次成分は44次元である。
この主張は剰余傾きと剰余体の拡大次数に制限を課さない。
\end{theorem}
\begin{proof}
局所計算の $K=(z^2+z^{-2})/2$, $U=K^2$ に対して傾きは
\[
 \Phi(z)=\frac{(z^8+14z^4+1)^3}{4z^4(z^4-1)^4}=\frac{P(z)}{Q(z)} .
\]
前節の十一桁の消去で単元として使用した定数項の分母と分子、および
選んだ二次小行列式の先頭係数を記録する。文字0,1の未切断六次射影と
飽和五次イデアルの積の階数8の主元も記録する。
\path{computations/verify_rank_fifteen_integral_slope.py} は整数係数の各多項式の
内容が2の冪であることを検査し、原始因子を重複なく合わせた多項式 $F$ を作る。
その次数は 205 である。記録した主元での消去と $H$ 除算前の厳密な零は
$\mathbb Z[1/2,z,1/F][[H]]$ で同じ証明を与える。従って $F(\bar z)\ne0$
となる表示が一つあれば十分である。

$i^2=-1$ とし、$z\mapsto iz$ と $z\mapsto(z+1)/(z-1)$ のいずれでも
$\Phi$ が変わらないことは式を代入すれば分かる。有限の傾きでは
$z(z^4-1)=0$ の根は現れないので、まずその因子を除く。
$f$ と $f(-z)$ の最大公約元を取れば、零点が $z,-z$ の双方に残る。
モニック化したこの多項式は偶多項式なので、$iz$ の代入は偶数次係数の符号だけで扱える。
さらにその代入および $(z+1)/(z-1)$ の斉次代入との最大公約元を取る。
次数が下がる限りこれを反復する。全ての変換先が悪い点は各段で残るため、
最終の多項式が単元なら必ず良い表示がある。

有理数上では最終の最大公約元は1である。整数多項式 $a,b$ の
最大公約元 $d$ ごとに $a,b$ の先頭係数、内容と
$\operatorname{Res}(a/d,b/d)$ の素因子を検査する。
これらが単元となる素数では剰余でも同じ消去が有効である。
残る奇素数はコードに列挙し、各々の素体の多項式環で最初から最大公約元を取り直す。
残った根は全て、以下の六つの $\bar D$ に対する $P-\bar D Q$ の根で尽きることを
反復する最大公約元で検証する。この検査は剰余体の代数閉包でも有効である。

最初の主元で残った $(p,\bar D)$ と、その傾きの追加表示を与える
$\bar z$ の最小多項式 $m$ は順に
\[
\begin{array}{c|cccccc}
p&5&7&13&19&29&37\\
\bar D&0&0&0&17&3&0\\
m&z^2+z+2&z^2+z+6&z-2&z-2&z-2&z-2
\end{array}
\]
である。各 $m$ の既約性、$m\mid P-\bar D Q$ と $m\nmid Q$ を検査する。
19以外の五つの剰余体で同時に可逆な主元を取り、19ではその表示で可逆な主元を取って、有理関数体上の十一桁の計算をやり直す。
中心方程式の自由列をこの主元に従って選び直し、文字零の後二自由列と文字一の両自由列を使う。
全十六個の六次射影で目標次元に達し、五次中心イデアルの積の最終消去で
必要な除算回数はどの始点でも文字順に $(0,0,2,2)$ である。

所定の有限な $D$ に対し $P-DQ$ はモニックで定数項1である。
付値を代数閉包に延長すれば全根は整な単元であり、積の還元から
任意の指定した剰余根が持ち上がる。奇標数で $z^4=1$ と $P=0$ は両立しない。
よって $U-1$ は単元で、$H^6=y(U-1)^4$ は正付値を与える。
$x\ne0$ なら $H^2=x(U-1)^2/(U+3)$ と取る。
$x=0$ は定理\ref{thm:rank-fifteen-x-zero}でも認証済みである。
前節の経路・中心の下界と選んだ四式の上界を合わせ、同じ正則性座標の窓へ移せば結論を得る。
\end{proof}
